\section{内陸地殻内地震}

\subsection{平成26年度 問32}
\label{q:h26-reg-applied-32}

\noindent\textbf{原典：} sources/H26/17\_applied\_science.pdf、PDF物理13頁。

\textbf{公式正答：} 5（single）。

\par\textbf{検討未確定：} 正答5の「最も活発」「約1000年」は母集団と平均期間の定義が不足する。公式5と独立初回5は保持し、厳密な下限の主張は未解決とする。

\paragraph{問題の要旨} 内陸地殻内地震の発生位置・規模・深さと、同じ活断層の再来間隔を区別する。

\Needspace{4\baselineskip}\paragraph{確認した条件・記述（独自要約）}\mbox{}\par

\noindent\textbf{確認範囲：} 問題幹、数値、単位、否定、全5候補を原典画像で照合した。\par

\noindent\textbf{成立条件：} 内陸という語を震央が必ず陸上という意味に限定しない。気象庁マグニチュードとモーメントマグニチュードを混ぜない。再来間隔は全国の地震間隔でなく一つの断層である。\par

\paragraph{問われる核心} 内陸地殻内地震の発生位置・規模・深さと、同じ活断層の再来間隔を区別する。

\paragraph{必要知識} 内陸という語を震央が必ず陸上という意味に限定しない。気象庁マグニチュードとモーメントマグニチュードを混ぜない。再来間隔は全国の地震間隔でなく一つの断層である。

\paragraph{考え方・正答を選ぶ理由（査読待ち初稿）} 地殻内地震は上盤側プレートの内部でも発生し、プレート境界地震と異なる。1995年兵庫県南部地震の気象庁マグニチュードは7.3であり、7を超えないという2には出題以前の反例がある。同じ活断層の活動間隔は通常人の寿命より長く、公式は5である。ただし「最も活発でも約1000年」を厳密な下限とするには断層群の定義と歴史資料の追加確認が必要である。

\Needspace{4\baselineskip}\paragraph{選択肢の検討（査読待ち初稿）}\mbox{}\par

\noindent\textbf{1：不適切} プレート境界で発生。地殻内の活断層活動は上盤側のプレート内部でも起きる。境界面のすべりとの区別が必要である。\par

\noindent\textbf{2：不適切} 最大級でもM7を超えない。出題前の兵庫県南部地震の気象庁M7.3が反例である。\par

\noindent\textbf{3：不適切} 震源はほとんど20 kmより深い。上部地殻の脆性変形を主体とする浅い地震であり、厚い地殻のすべてが地震を起こすわけではない。\par

\noindent\textbf{4：不適切} 震央は陸上。地殻内地震という発生機構の区分には沿岸の海域に震央を持つものも含み得る。地表の海陸区分だけで機構を決めない。\par

\noindent\textbf{5：解釈要確認} 最も活発な活断層でも約1000年に1回。公式が採る代表的な長期尺度である。同一断層の繰返しと各地の地震を区別する意図は妥当だが、全活断層の下限が1000年であるとの定量保証は未確定である。\par

\paragraph{再利用する考え方} 発生機構・震源の深さ・震央の海陸・Mの尺度を分ける。平均間隔が長いことは近日の発生が不可能という意味ではない。

\paragraph{共通理論} \hyperref[theory:interior:crust_quakes]{第\ref*{theory:interior:crust_quakes}節}

\par 参考：\cite{jma_hyogo_history,herp_longterm_intervals}。

\par 難易度 2（校準初稿）、優先度未評価。
